% ============================================================================
% Fichier  : partie6/P06_C26_mesurer.tex
% Titre    : Mesurer la Sécurité
% Auteur   : MINKA MI NGUIDJOÏ Thierry Emmanuel
% Version  : 1.0 -- novembre 2026
% ============================================================================
% Sources : ISO/IEC 27004:2016 ; NIST SP 800-55 Rev. 1 (2008) ; instrument
% OUI / NON / NEK et mesures mu, U, H (Minka Mi Nguidjoi, 2026, A1, HAL) ;
% programme de l'auditologie (A0, HAL). Seuls ces deux textes de l'auteur,
% déposés, sont cités. Valeurs de mu, U, H et intervalles de Wilson
% calculés le 07/11/2026.
% ============================================================================

\chapter{Mesurer la Sécurité}
\label{chap:mesurer}
\chaptermark{Mesurer la sécurité}

\epigraph{%
    « Ce que l'on ne sait pas n'est ni conforme, ni non conforme~:\\
    c'est inconnu, et il faut le dire. »%
}{Principe du cours}

\niveauChapitre{Expert}{Chapitres~\ref{chap:gouvernance},
\ref{chap:normes} et~\ref{chap:evaluer}.}

\begin{boxobjectifs}
\begin{itemize}
    \item Construire des indicateurs de sécurité selon l'ISO/IEC~27004.
    \item Distinguer indicateurs de mise en œuvre, d'efficacité et
          d'impact.
    \item Utiliser l'instrument \OUI{}~/ \NON{}~/ \NEK{} pour mesurer ce
          que l'on sait d'un dispositif.
    \item Calculer et interpréter la conformité \(\mu\), l'incertitude
          \(U\) et l'entropie \(H\).
    \item Dire avec quelle précision une mesure est connue.
\end{itemize}
\end{boxobjectifs}

\bigskip

% ============================================================================
\section*{Question directrice : LiZbiZ est-elle mieux protégée six mois après~?}
% ============================================================================

En septembre, le conseil de LiZbiZ demande si les mesures prises depuis
mars ont amélioré la sécurité. Le RSSI pourrait répondre par un
pourcentage de mesures « en place ». Mais une partie des réponses
recueillies auprès des services sont des « je ne sais pas »~: faut-il les
compter comme des oui, comme des non, ou les écarter~? Ce chapitre
répond, et clôt le cours.

\medskip
\textbf{Plan du chapitre.}
\begin{enumerate}
    \item Les indicateurs de sécurité (section~\ref{sec:c26-indicateurs}).
    \item Le problème de l'ignorance (section~\ref{sec:c26-ignorance}).
    \item Trois mesures~: \(\mu\), \(U\), \(H\) (section~\ref{sec:c26-mesures}).
    \item La précision d'une mesure (section~\ref{sec:c26-precision}).
    \item Synthèse du cours (section~\ref{sec:c26-synthese}).
\end{enumerate}

\separateur

% ============================================================================
\section{Les Indicateurs de Sécurité}
\label{sec:c26-indicateurs}
% ============================================================================

L'ISO/IEC~27004 et le NIST SP~800-55 distinguent plusieurs familles
d'indicateurs~\cite{iso27004_2016,nist80055}.

\begin{tableaucours}{L{3.4cm}L{4.6cm}L{4.2cm}}{%
    \tableauHeader{Famille} &
    \tableauHeader{Question} &
    \tableauHeader{Exemple}
}
Mise en œuvre & La mesure existe-t-elle~? & Part des comptes distants
avec second facteur \\
Efficacité & Fonctionne-t-elle~? & Part des restaurations réussies dans
le RTO \\
Impact & Change-t-elle le résultat~? & Délai de détection des incidents \\
\end{tableaucours}
\vspace{-0.8em}
\begin{center}{\small\itshape Tableau~26.1. Familles d'indicateurs}\end{center}

Un indicateur de mise en œuvre est facile à mesurer et peu informatif~;
un indicateur d'impact est informatif et difficile. Un tableau de bord
équilibré combine les trois (Chapitre~\ref{chap:gouvernance}).

% ============================================================================
\section{Le Problème de l'Ignorance}
\label{sec:c26-ignorance}
% ============================================================================

Quand on évalue un dispositif par des questions (« les sauvegardes
sont-elles testées~? »), certaines réponses ne peuvent pas être établies~:
la pièce n'existe pas, le prestataire ne répond pas, le journal a été
effacé. Les grilles habituelles ne laissent le choix qu'entre oui et non~;
elles obligent alors à deviner, et le plus souvent à répondre oui.

\begin{boxdef}
\termecle{Instrument trichotomique}~: grille dont chaque question admet
trois réponses~: \OUI{} (la mesure est en place, preuve à l'appui),
\NON{} (elle ne l'est pas) et \NEK{} (\emph{not enough knowledge}~:
connaissance insuffisante pour conclure). Il a été formalisé par l'auteur
dans une prépublication déposée sur HAL~\cite{minka2026nek}, dans le cadre
d'un programme de recherche plus large~\cite{minka2026auditologie}.
\end{boxdef}

% ============================================================================
\section{Trois Mesures : \texorpdfstring{\(\mu\), \(U\), \(H\)}{mu, U, H}}
\label{sec:c26-mesures}
% ============================================================================

Sur \(n\) questions, avec \(n_O\) réponses \OUI{}, \(n_N\) \NON{} et
\(n_K\) \NEK{}~:

\begin{itemize}
    \item \(\mu = n_O / n\)~: la \termecle{conformité}, part des réponses
          \OUI{}~;
    \item \(U = n_K / n\)~: l'\termecle{incertitude}, part des réponses
          \NEK{}~;
    \item \(H = -\sum p \log_2 p\), sur les trois proportions~:
          l'\termecle{entropie}, en bits, au plus \(\log_2 3 \approx 1{,}585\),
          qui mesure la dispersion des réponses.
\end{itemize}

\begin{boxlecture}{Les mesures de sécurité de LiZbiZ, mars et septembre}
Vingt questions sur les mesures clés (second facteur, correctifs,
sauvegardes, journaux, comptes, contrat du prestataire\ldots), posées
aux mêmes services.

\medskip
\begin{tableaucours}{L{2.6cm}cccc}{%
    \tableauHeader{Date} &
    \tableauHeader{O / N / K} &
    \tableauHeader{\(\mu\)} &
    \tableauHeader{\(U\)} &
    \tableauHeader{\(H\) (bit)}
}
Mars & 4 / 6 / 10 & 0,20 & 0,50 & 1,485 \\
Septembre & 12 / 3 / 5 & 0,60 & 0,25 & 1,353 \\
\end{tableaucours}

\medskip
\textbf{Constatation.} La conformité passe de 0,20 à 0,60 et
l'incertitude de 0,50 à 0,25. \textbf{Analyse} (confiance moyenne). La
progression est réelle, mais un quart des réponses reste inconnu~:
présenter « 60~\% de mesures en place » sans dire que 25~\% sont
inconnues donnerait au conseil une image trop favorable. Compter les
\NEK{} comme des \OUI{} afficherait 85~\%~; les compter comme des \NON{},
60~\% aussi, mais en cachant qu'on ne sait pas.

\hfill\lecture{0{,}60}{0{,}25}{1{,}35}
\end{boxlecture}

\begin{boxinfo}
\textbf{Lire les trois mesures ensemble.} \(U\) dit s'il faut d'abord
chercher la preuve~; \(\mu\) dit ce qui est en place parmi ce que l'on
sait~; \(H\) dit si les réponses sont homogènes ou dispersées. Une
réduction de \(U\) est un progrès de la connaissance, même si elle fait
apparaître des \NON{}.
\end{boxinfo}

% ============================================================================
\section{La Précision d'une Mesure}
\label{sec:c26-precision}
% ============================================================================

Vingt réponses, c'est peu. Une proportion observée sur un petit nombre de
cas est connue avec une marge large. L'intervalle de Wilson, adapté aux
petits effectifs, le montre~: pour \(U = 0{,}50\) sur 20~réponses,
l'intervalle à 95~\% va environ de 0,30 à 0,70~; pour \(U = 0{,}25\), de
0,11 à 0,47. Les deux intervalles se recouvrent~: la baisse de
l'incertitude est probable, mais 20~questions ne suffisent pas à
l'établir avec certitude. Un rapport rigoureux le dit, et recommande
d'élargir l'échantillon pour les domaines importants.

\begin{boxalert}
\textbf{Violation typique.} Présenter « 60~\% » au conseil comme une
mesure exacte.

\smallskip
\textbf{Règle~:} une mesure s'accompagne de son effectif et, quand il est
petit, de sa marge. Les seuils qui servent à classer (« acceptable à
partir de 80~\% ») sont des conventions, à déclarer comme telles.
\end{boxalert}

% ============================================================================
\section{Synthèse du Cours}
\label{sec:c26-synthese}
% ============================================================================

\begin{tableaucours}{L{2.8cm}L{9.4cm}}{%
    \tableauHeader{Partie} &
    \tableauHeader{Ce qu'elle a apporté face à l'Opération GAROUA}
}
I. Fondements & Nommer les pertes (propriétés), comprendre la chaîne
d'attaque, savoir ce que la cryptographie protège ou non \\
II. Gouvernance et risques & Un responsable, une politique, une
appréciation des risques qui décrivait l'attaque d'avance, des actifs
connus, des personnes formées \\
III. Normes et droit & Des référentiels pour ne rien oublier~; des
obligations (journaux, confidentialité) et des infractions \\
IV. Sécurité technique & Couper chaque maillon~: second facteur,
segmentation, correctifs, comptes, détection, réponse \\
V. Défis avancés & Reprendre (sauvegardes immuables, RTO), préserver la
preuve, juger les technologies nouvelles \\
VI. Mesurer & Savoir si l'on progresse, et dire ce que l'on ne sait pas \\
\end{tableaucours}
\vspace{-0.8em}
\begin{center}{\small\itshape Tableau~26.2. Les six parties du cours et
l'Opération GAROUA}\end{center}

\begin{boxmetier}{RSSI}
\textbf{Rendre compte honnêtement.} Le RSSI présente au conseil la
conformité, l'incertitude et leur évolution, avec l'effectif et la
marge, et un plan pour réduire d'abord les \NEK{} sur les mesures
critiques.
\end{boxmetier}

\begin{boxmetier}{Auditeur de sécurité}
\textbf{Coter \NEK{} quand il le faut.} L'auditeur cote \NEK{} une mesure
dont il n'a pas pu obtenir la preuve, au lieu de la supposer en place~;
il indique ce qu'il faudrait pour conclure.
\end{boxmetier}

\separateur

\begin{boxresume}
\textbf{Ce qu'il faut retenir de ce chapitre~:}
\begin{itemize}
    \item \textbf{Les indicateurs} mesurent la mise en œuvre, l'efficacité
          et l'impact~; un tableau de bord les combine.
    \item \textbf{L'instrument \OUI{}~/ \NON{}~/ \NEK{}} empêche de
          deviner quand on ne sait pas.
    \item \textbf{\(\mu\), \(U\) et \(H\)} se lisent ensemble~; réduire
          \(U\) est un progrès de la connaissance.
    \item \textbf{Une mesure} s'accompagne de son effectif et de sa marge~;
          les seuils sont des conventions.
    \item \textbf{La sécurité} est un cycle~: connaître, protéger,
          détecter, répondre, reprendre, mesurer.
\end{itemize}
\end{boxresume}

\begin{boxinfo}
\textbf{Points de contrôle de l'élève.} Je sais~: distinguer les trois
familles d'indicateurs~; coter \OUI{}, \NON{} ou \NEK{}~; calculer
\(\mu\), \(U\) et \(H\)~; expliquer pourquoi compter les \NEK{} comme des
\OUI{} est trompeur~; dire pourquoi 20~réponses donnent une mesure
imprécise.
\end{boxinfo}

\bigskip

\begin{boxexo}[title={\ding{45}\quad Exercice~26.1. Quelle famille~? \niveaubase}]
Classez~: nombre de postes avec un EDR~; délai médian de détection~; part
des alertes traitées en moins de quatre heures.
\end{boxexo}

\begin{boxexo}[title={\ding{45}\quad Exercice~26.2. Coter \niveaubase}]
Cotez \OUI{}, \NON{} ou \NEK{}~: (a) la politique de sauvegarde existe,
mais aucun test de restauration n'a été fourni~; (b) le prestataire
déclare avoir activé le second facteur, sans preuve~; (c) l'extraction de
l'annuaire montre trois comptes orphelins.
\end{boxexo}

\begin{boxexo}[title={\ding{45}\quad Exercice~26.3. Calculer \niveaubase}]
Calculez \(\mu\), \(U\) et \(H\) pour 14~\OUI{}, 4~\NON{} et 2~\NEK{}.
\end{boxexo}

\begin{boxexo}[title={\ding{45}\quad Exercice~26.4. Les cas limites \niveaumoyen}]
Que valent \(\mu\), \(U\) et \(H\) quand toutes les réponses sont \NEK{}~?
Pourquoi une entropie nulle ne signifie-t-elle pas ici une certitude~?
\end{boxexo}

\begin{boxexo}[title={\ding{45}\quad Exercice~26.5. Compter les NEK \niveaumoyen}]
Reprenez les données de septembre. Calculez la conformité affichée si
l'on compte les \NEK{} comme des \OUI{}, puis comme des \NON{}. Quelle
présentation recommandez-vous au conseil~?
\end{boxexo}

\begin{boxexo}[title={\ding{45}\quad Exercice~26.6. La marge \niveaumoyen}]
Expliquez à M.~ETEME, sans formule, pourquoi 20~réponses ne suffisent pas
à affirmer que l'incertitude a diminué.
\end{boxexo}

\begin{boxexo}[title={\ding{45}\quad Exercice~26.7. Un indicateur d'impact \niveaumoyen}]
Proposez un indicateur d'impact pour la mesure « double authentification
de l'accès distant », et la manière de le mesurer.
\end{boxexo}

\begin{boxexo}[title={\ding{45}\quad Exercice~26.8. Réduire U \niveaumoyen}]
Pour les cinq \NEK{} de septembre, proposez la preuve qui permettrait de
conclure, et classez-les par priorité.
\end{boxexo}

\begin{boxexo}[title={\ding{45}\quad Exercice~26.9. Le tableau de bord \niveauexpert}]
Construisez le tableau de bord trimestriel du conseil de LiZbiZ~: six
indicateurs des trois familles, et la lecture \(\mu\), \(U\), \(H\) des
vingt mesures clés.
\end{boxexo}

\begin{boxexo}[title={\ding{45}\quad Exercice~26.10. Synthèse \niveauexpert}]
Rédigez, en deux pages, le rapport de clôture de l'Opération GAROUA~: ce
qui s'est passé, ce qui a été corrigé, ce qui reste inconnu, et les
priorités de l'année suivante.
\end{boxexo}

\bigskip

\begin{boxinfo}
\textbf{Pour aller plus loin.}
\begin{itemize}
    \item ISO/IEC~27004 et NIST SP~800-55~\cite{iso27004_2016,nist80055}.
    \item L'instrument \OUI{}~/ \NON{}~/ \NEK{}~\cite{minka2026nek} et le
          programme de l'auditologie~\cite{minka2026auditologie}.
    \item Les UE INF3226 (audit des systèmes d'information) et SEC4031
          (investigation numérique), qui prolongent ce cours.
\end{itemize}
\end{boxinfo}

\begin{boxinfo}
\textbf{Conclusion du cours.} Sécuriser un système d'information, c'est
connaître ce que l'on protège, contre qui, avec quelles mesures, et savoir
vérifier qu'elles fonctionnent. Les annexes rassemblent le glossaire, les
outils, les modèles et les textes utiles à la pratique.
\end{boxinfo}
